目标区域为半径 $1800\unit{m}$ 的圆域，干扰源频道互异。机器狗以 $5\unit{m/s}$ 直线移动，在静止点逐频道检测；一次频道切换、检测和成功清除分别耗时 $1\unit{s}$、$5\unit{s}$ 和 $5\unit{s}$。测得示向度与真实方位角的误差属于 $[-1^\circ,1^\circ]$，同一地点的误差固定。干扰源有效接收半径未知但在 $[1000,1500]\unit{m}$ 内。

四问形成递进关系。第一问要求由多次示向观测构造定位区域并计算直径，同时判断半径为直径一半的圆是否必然覆盖。第二问在只有一次全向源示向观测时选择第二检测点，需要同时控制收不到信号的风险和两条方位线接近平行的病态性。第三问要求在源数未知但属于 $10$--$16$ 的条件下，对全向源给出保证无遗漏的有限搜索、定位与清除算法。第四问加入辐射半平面方向未知的定向源，要求重新证明发现的完备性，且不能把无信号简单解释为附近无源。

对每一问分别交付几何算法或机器狗策略、可复算数值结果和验证证据。问题三、四以清除比例为首要约束，在全部清除的前提下比较虚拟总时间与平均定位清除时间。
